\documentclass[11pt,a4paper]{article}
\usepackage[utf8]{inputenc}
\usepackage[T1]{fontenc}
\usepackage{lmodern}
\usepackage[margin=2.4cm]{geometry}
\usepackage{graphicx}
\usepackage{booktabs}
\usepackage{tabularx}
\usepackage{amsmath,amssymb}
\usepackage[round,authoryear]{natbib}
\usepackage{xcolor}
\usepackage[hidelinks]{hyperref}
\usepackage{caption}
\usepackage{microtype}
\usepackage{enumitem}
\usepackage[section]{placeins}
\captionsetup{font=small,labelfont=bf,labelsep=period}
\setlist{nosep,leftmargin=1.4em}
\renewcommand{\arraystretch}{1.15}
\newcommand{\figpath}{figs}
\newcommand{\x}{\ensuremath{\times}}

\title{\textbf{LandEmu v1: a learned emulator of the trait-based,\\ individual-tree vegetation model LPJmL-FIT}\\[4pt]\large Daily carbon and water exchange and forest structure at a fraction of the cost}
\author{Jamir Priesner\\[2pt]
\small Potsdam Institute for Climate Impact Research (PIK), Member of the Leibniz Association,\\
\small Potsdam, Germany\\
\small Correspondence: jamir.priesner@chance-international.org}
\date{Manuscript draft, 21 October 2026}

\begin{document}
\maketitle

\begin{abstract}
\noindent
Individual-based, trait-based vegetation models resolve how forests of differently adapted trees grow, compete and
die, but they are too expensive to run inside Earth system models. We present LandEmu v1, a learned emulator of the
dynamic global vegetation model LPJmL-FIT that reproduces both its daily carbon and water exchange and its yearly
forest structure and trait composition, driven by daily weather alone. Its design follows from a systematic review of
115 emulator and hybrid-model studies of land models and from measurements that ruled out an earlier hybrid design:
re-implementing the original's daily physics made it 4.6 times slower, and no single process of the original takes
enough run time to give even a twofold speed-up. A five-member ensemble of small neural networks therefore replaces the
entire daily step, with snow, top-metre and deep soil water carried as explicit stores that close the water balance
exactly. Once a year, gradient-boosted tree models predict tree density, biomass per tree, four trait medians, leaf
area, vegetation carbon and plant-type cover as a departure from a per-cell anchor, driven by the trailing 30-year
climate; this part is not recursive, so its errors cannot accumulate. Tested on a climate model and an independent run
of LPJmL-FIT withheld from their training, the yearly maps' error in dense forests is 1.16 (tree density) and 1.07
(biomass per tree) times that of a further run of the stochastic original, and smaller for the trait medians. Tested on
cells withheld from their training, under climate models they had seen, the daily networks' 30-year mean fluxes are 2.4
(gross primary production) to 5.2 (runoff) times further from the original than a further run is --- about 3--6\,\%
per cell for production and evapotranspiration, more for runoff --- with area totals within 2.5\,\%. After the start
state has faded the structure error grows only slowly, with a creep in tree density under strong warming, and the flux
error grows only for gross production under the two warmest climates. For a nitrogen-limited configuration of
LPJmL-FIT with rising CO$_2$, the emulator reproduces 81--97\,\% of the change in vegetation carbon under two unseen
climate models (a second run of the original: 61--89\,\%), but damps the largest warming-only loss (57--72\,\% of it)
and the cell-to-cell spread of the response; the pre-registered test was not passed. The emulator costs
$7.6\times10^{-3}$ core-seconds per cell-year, about 35 and 900 times less than LPJmL-FIT at 25 and 500 patches.
\end{abstract}

\section{Introduction}\label{sec:intro}

Terrestrial ecosystems take up about a third of anthropogenic CO$_2$ emissions \citep{Friedlingstein2023GCB}, and
the vegetation that does so is not a fixed surface but a population of plants that grow, compete, die and are
replaced. Dynamic global vegetation models (DGVMs) represent this population in Earth system models (ESMs), usually
through a few plant functional types with fixed properties. A growing body of work argues that this is not enough:
the size structure and the functional diversity of forests shape how they respond to drought and warming, and
models that resolve individuals, cohorts and traits behave differently from those that do not
\citep{Fisher2018,Fisher2020}. LPJmL-FIT, a version of the LPJmL model with individual trees and variable traits
\citep{Sakschewski2015,Sakschewski2016,Thonicke2020}, is one of these models: each tree carries its own specific leaf
area, wood density, rooting depth and drought threshold, drawn when it establishes, and the forest's trait
composition shifts as trees with poorly suited traits die. Amazon forests with more trait diversity recover better
from drought in this model than forests without it \citep{Sakschewski2016}.

The price of this realism is cost. LPJmL-FIT simulates every tree every day, and because establishment and mortality
are random it averages over many replicate patches per grid cell: 25 in the runs used here, around 500 in published
applications. Its run time grows in proportion to the number of patches, which puts it orders of magnitude beyond
what an ESM can afford for its land surface. Emulation --- replacing the expensive model with a fast statistical or
machine-learned approximation --- is the obvious route, and the land-modelling literature contains many emulators.
A systematic review of 115 emulator and hybrid studies of DGVMs and land surface models \citep{PriesnerReview}
shows, however, that most of them serve calibration and sensitivity analysis: 89 emulate only one or two variables,
a single domain or a single site, or a summary statistic. Only seven emulate a land model broadly enough to stand in
for part of it, and each covers only part of the model: prognostic soil states and turbulent fluxes of a
weather-prediction land scheme \citep{R00067,R00593}, carbon and stand structure of a demographic model
\citep{R01133,R00514}, biome-level carbon for integrated assessment \citep{R00661}, equilibrium carbon pools for
spin-up \citep{R01211} or steady-state parameter responses \citep{R00004}. None emulates size or trait distributions
or individual mortality; the closest is a state-transition network for the landscape model iLand
\citep{Rammer2019}. None of the seven has been coupled to an ESM, only 12 of the 115 studies test climates outside the
training range, and only 25 check conservation.

The review also distils what goes wrong, and these lessons shaped the emulator described here. Tree ensembles trained
to predict one step ahead drift when run freely at long lead times (up to about a year in the cited tests), while
networks trained on their own multi-step rollouts do not \citep{R00067,R00593}. Agreement in the present climate says little about agreement under a changed
climate: an emulator of the ecosystem model ecosys matched present-day net primary production almost exactly yet
underestimated it by 204\,gC\,m$^{-2}$\,yr$^{-1}$ by 2100 \citep{R00289}, and an LPJ-GUESS emulator trained on
stylised climate perturbations failed to transfer to realistic scenarios \citep{R00430}. Offline skill can collapse
once feedbacks act, as for a machine-learned fire model in a land surface model whose spatial correlation fell from
0.96--0.98 to 0.59--0.70 \citep{R00911,R00342}. Water and carbon budgets should be closed by construction rather than
encouraged by a loss term \citep{R00755}. From this evidence the review derives six requirements for emulators that
are to be used inside ESMs: prognostic state and memory, closure of the coupling interface, conservation, stability
under feedback, credibility under out-of-distribution climate, and quantified uncertainty with usable software
interfaces.

This paper describes LandEmu v1, an emulator of LPJmL-FIT built against these requirements. Its contributions are:
\begin{enumerate}
  \item a complete emulator of an individual-tree, trait-based DGVM (Fig.~\ref{fig:overview}) --- daily gross and net primary production,
        interception, transpiration, soil evaporation, runoff, snow and soil water, together with yearly tree density,
        biomass per tree, four trait medians, leaf area, vegetation carbon and plant-type cover --- that runs from
        daily weather alone and conserves water exactly;
  \item an account of how its design followed from measurements that ruled out a hybrid design (Sect.~\ref{sec:path});
  \item an evaluation that holds out a whole climate model and an independent run of the original (yearly maps) or
        whole blocks of cells (daily networks), and that judges every error against the original's own run-to-run
        variability (Sect.~\ref{sec:eval});
  \item a test of whether the emulator learns the warming response when the training data cover the test climate,
        and a full-grid application (Sect.~\ref{sec:results});
  \item an end-to-end cost measurement on one core against the original at its training and publication settings.
\end{enumerate}

\begin{figure}[!t]
  \centering
  \includegraphics[width=\textwidth]{\figpath/fig01_overview.pdf}
  \caption{How LandEmu advances one year. (A) The daily water--carbon model: the inputs of day $t$ pass through an
  ensemble of five neural networks whose bounded outputs drive a closed three-store water bucket; the stores are
  carried to the next day. (B) The yearly forest maps: 25 climate statistics of the year enter a trailing 30-year
  window, and 18 gradient-boosted tree models predict the forest structure and stand as a departure from a per-cell
  anchor. The stand (leaf area, vegetation carbon, cover of ten plant types) is the daily model's input in the
  following year (dashed). Every other quantity is predicted afresh each year.}
  \label{fig:overview}
\end{figure}

\section{The model to be emulated: LPJmL-FIT}\label{sec:target}

LPJmL-FIT couples the water, carbon and vegetation dynamics of LPJmL \citep{Schaphoff2018} with an individual-tree
forest module \citep{Sakschewski2015,Thonicke2020}. On each 0.5$^\circ$ grid cell the forest grows on a set of
replicate patches (225\,m$^2$ each). Trees of seven woody plant types establish with traits drawn within type-specific
ranges, or inherited with small variations from large trees of the same patch; they compete for light in a layered
canopy and for water in a 23-layer soil, allocate carbon each year, and die from carbon starvation, age, heat and
drought stress. Three grass types fill the remaining space. The daily step computes photosynthesis and stomatal
conductance per tree, interception, transpiration, soil evaporation, snow, infiltration, percolation and runoff; the
yearly step allocates carbon, applies mortality and establishes new trees. We use the configuration with natural
vegetation only. For most of this paper nitrogen limitation is switched off and atmospheric CO$_2$ is held constant
in the future runs (at its 2019 level; 2014 in the runs of Sect.~\ref{sec:global}); this is how LPJmL-FIT is usually run, because without nitrogen limitation its CO$_2$
fertilisation is unbounded. The emulator of this configuration therefore takes no CO$_2$ input, which is faithful to
the original rather than a gap. Section~\ref{sec:resp} uses a second configuration, with nitrogen limitation
\citep{vonBloh2018} and rising CO$_2$, for which the emulator receives CO$_2$ as an input.

\paragraph{Where the run time goes.} We profiled the unmodified production binary at five contrasting cells
(temperate, tropical, semi-arid, Mediterranean, boreal; 25 patches, one core) and attributed the exclusive run time
to processes, assigning the statically linked mathematical library calls (21--27\,\% of all time) to the routines that
call them (Fig.~\ref{fig:runtime}a). Leaf gas exchange takes 36--46\,\% of the run time and soil water another
15--27\,\%. The yearly demography --- allocation, mortality, establishment and turnover --- takes only 0.4--1.1\,\%,
and the daily loop as a whole 97--98\,\%. The marginal cost is 0.20--0.33 core-seconds per cell-year at 25 patches,
and it grows linearly with the number of patches. Even if the eight most expensive processes were free of cost, the
model would speed up at most 30-fold at the temperate cell (Fig.~\ref{fig:runtime}b), whereas an intermediate-resolution atmosphere
leaves an allowance of about 0.0135--0.030 core-seconds per land cell-year (10\,\% of a measured coupled
atmosphere's cost at T31 and T63 resolution). Removing the 25-patch ensemble alone is worth a factor of 25, and of
about 500 at the publication setting.

\paragraph{The original is stochastic.} Two independent runs of LPJmL-FIT --- each with its own 1000-year spin-up and
random seed, under identical climate --- disagree on a typical cell by about 6\,\% in the number of trees taller than
5\,m, 12\,\% in biomass per tree, 2--10\,\% in the trait medians and 1.0--1.6\,\% in 30-year mean fluxes. In sparse
stands the disagreement is far larger. Any emulator can at best predict the expectation over such runs; this sets
both the target and the yardstick of the evaluation (Sect.~\ref{sec:eval}).

\begin{figure}[!tbp]
  \centering
  \includegraphics[width=\textwidth]{\figpath/fig02_runtime.pdf}
  \caption{Where LPJmL-FIT spends its time. (a) Exclusive run time by process at five biome cells (25 patches, one
  core; mathematical library calls attributed to their callers). (b) The largest speed-up obtainable at the temperate
  cell if processes were made free of cost one after another, against the speed-up needed for an intermediate-resolution
  (T31) atmosphere, the gain from dropping the 25-patch ensemble, and the speed-up needed at 500 patches.}
  \label{fig:runtime}
\end{figure}

\section{From a hybrid to a fully learned emulator}\label{sec:path}

LandEmu is the second architecture we built for LPJmL-FIT. The first was a hybrid, the design favoured by much of the
hybrid-modelling literature \citep{Reichstein2019,Shen2023}: learn the slow, poorly constrained part and keep the
physics. It had three components (Fig.~\ref{fig:path}, left). A learned yearly demography predicted the number of
trees per patch with a distributional random forest, sampled new trees' traits from a Gaussian copula, and applied the
original's mortality rules, ported from its source code. The daily physics --- photosynthesis, the stomatal solve, the layered
canopy, soil water, phenology, allocation and growth --- was re-implemented as differentiable code, validated
process by process against the original at five biome cells. An energy closure supplied the skin temperature and
sensible heat that LPJmL-FIT does not compute. The training data were two runs of the original under one historical
and one future climate. Five findings --- four measurements and one from the literature --- ruled this design out
(Fig.~\ref{fig:path}, middle).

\emph{(i) It was not faster.} The differentiable re-implementation cost 1.23 core-seconds per cell-year at 25 patches,
4.6 times more than the original. Per tree, its daily step cost 51--58 times the original's (83\,\% of it in the
iterative stomatal solve with a finite-difference derivative); only because it carried an order of magnitude fewer
trees (those taller than 5\,m) was the overall factor as small as 4.6. Measured solver improvements would have reached
parity at best.

\emph{(ii) The original's cost is the daily step and the patch ensemble} (Sect.~\ref{sec:target}). Any substantial
speed-up must replace the daily step as a whole and must predict the expectation over patches instead of simulating
replicates.

\emph{(iii) One-step learning drifts when run freely.} The tree-count model reached an out-of-sample coefficient of
determination of 0.982 one year ahead --- but carrying the previous year's count forward already reached 0.962. Fed
its own output, its area-weighted warming response turned from +0.71 to $-0.23$ of the original's within decades,
because it followed 96\,\% of large increases but only 87\,\% of large decreases: a small, state-dependent asymmetry
that accumulates. This is the failure the review documents for one-step tree ensembles \citep{R00067}.

\emph{(iv) One climate per place cannot teach a response.} Across the global grid, 76\,\% of a cell's warming is
predictable from its present climate, so with one historical and one future run the climate response is confounded
with geography --- the identification problem behind the out-of-distribution failures in the literature
\citep{R00289,R00430}.

\emph{(v) The literature points the same way}: conservation by construction, out-of-range climate tests and explicit
state are what the published emulators most often lack \citep{PriesnerReview}.

From these findings follow five design principles (Fig.~\ref{fig:path}, right). \textbf{P1}: learn the whole daily
step with one small network per cell and day, keeping water stores explicit and closed. \textbf{P2}: predict the
expectation over the patch ensemble, so that the number of patches drops out of the cost. \textbf{P3}: do not make the
forest state recursive; predict each year's state as a departure from a per-cell anchor, driven by climate, so that
errors cannot accumulate. \textbf{P4}: train on climate contrast --- many climate models and scenarios and several
independent runs --- and test on a climate model and a run never seen. \textbf{P5}: measure every error against the
original's own run-to-run variability.

\begin{figure}[!tbp]
  \centering
  \includegraphics[width=\textwidth]{\figpath/fig_designpath.pdf}
  \caption{How the design was arrived at. Left: the components of the first, hybrid emulator. Middle: five findings
  that ruled it out, each a measurement on LPJmL-FIT or on the hybrid, or a result of the literature review. Right: the
  design principles of LandEmu that answer them.}
  \label{fig:path}
\end{figure}

\section{Emulator description}\label{sec:model}

LandEmu advances one year at a time in two blocks (Fig.~\ref{fig:overview}; Table~\ref{tab:io}). Block A, the daily
water--carbon model, runs 365 daily steps from the day's weather and last year's stand. Block B, the yearly forest
maps, runs once at the end of the year from climate statistics of the same weather and predicts the forest structure
and the stand, which becomes the daily model's input for the next year. Everything is driven by daily weather on the
0.5$^\circ$ grid: air temperature, precipitation, downward shortwave radiation and net longwave radiation; specific
humidity enters only the yearly climate statistics. All years have 365 days.



\begin{table}[!tbp]
\centering
\caption{What LandEmu reads, carries and returns.}\label{tab:io}
\small
\begin{tabularx}{\textwidth}{@{}lX@{}}
\toprule
\textbf{Daily input} & air temperature, precipitation, downward shortwave, net longwave (all days); specific humidity
(yearly statistics only); CO$_2$ in the nitrogen-on variant \\
\textbf{Per-cell set-up} & latitude, longitude, soil code, top-metre plant-available water capacity; anchors of every
yearly quantity; the climate of the anchor runs; the run's own 2000--2019 state and its offset from the anchor; initial
stores, stand and weather windows \\
\textbf{State carried} & snow, top-metre water, deep water; trailing weather sums and means (10, 30, 90, 365 days);
30-year window of yearly climate statistics; last year's stand \\
\textbf{Daily output} & gross primary production (GPP), net primary production (NPP), interception, transpiration, soil
evaporation, runoff, snow water equivalent, top-metre soil water, deep store \\
\textbf{Yearly output} & trees taller than 5\,m per patch, above-ground biomass per tree, medians of specific leaf area,
wood density, maximum rooting depth and drought threshold; stand leaf area index, vegetation carbon, fractional cover
of ten plant types \\
\bottomrule
\end{tabularx}
\end{table}

\subsection{Daily water--carbon model}\label{sec:daily}

\paragraph{Inputs.} For cell $c$ and day $t$ the network sees today's weather; trailing means of temperature (10 and
30 days) and shortwave radiation (30 days) and trailing sums of precipitation (30, 90 and 365 days), which stand in for
the memory of the deep soil without carrying it as a state; the day of year (as sine and cosine), the sine of latitude
and the day length; yesterday's snow $S_{t-1}$ and top-metre water $W_{t-1}$, $W_{t-1}$ divided by the capacity, and
the capacity itself; and the stand at the end of the previous year (leaf area index, vegetation carbon and the cover of
the ten plant types). No input is taken from the fluxes of the run being emulated. The top-metre water capacity, which
the original recomputes from its evolving soil carbon, is read from the original each year in the evaluation runs; it
changes over the century by a median 0.1--0.2\,\% (at most 4--6\,\% in single cells), so a fixed per-cell value would
serve.

\paragraph{Network.} Each of five ensemble members is a multilayer perceptron with three hidden layers of 256 units
and SiLU activations; all outputs share one trunk. The outputs are mapped to bounded rates: an interception fraction
$i\in[0,1]$, a snowfall fraction $f\in[0,1]$ and a melt fraction $m\in[0,1]$, non-negative transpiration $E_T$, soil
evaporation $E_S$ and runoff $R$, a signed exchange $Q$ between the top metre and the deep soil, non-negative GPP and
plant respiration $R_a$, so that $\mathrm{NPP} = \mathrm{GPP} - R_a$. The ensemble mean of the five members is used.

\paragraph{Closed water bucket.} With precipitation $P_t$, the stores are updated as
\begin{align}
  S_t &= S_{t-1}(1-m) + f P_t, \qquad E_{I,t} = i P_t, \\
  W_t &= W_{t-1} + P_t - (S_t - S_{t-1}) - E_{I,t} - E_{T,t} - E_{S,t} - R_t - Q_t, \\
  D_t &= \left(D_{t-1} + Q_t\right)\left(1 - 1/\tau\right) ,
\end{align}
where $D$ is a deep store that drains the amount $(D_{t-1}+Q_t)/\tau$ into runoff each day, with a time constant
$\tau = 365$\,days. Water above the top-metre capacity is added to runoff and a deficit below zero is taken from the
deep store. Hence, every day and to rounding precision,
\begin{equation}
  P_t = \Delta(S + W + D)_t + E_{I,t} + E_{T,t} + E_{S,t} + R^{\mathrm{tot}}_t ,
\end{equation}
with $R^{\mathrm{tot}}$ the total runoff. The bounds restrict nothing the original does: on eight million cell-days
interception and snow accumulation never exceeded precipitation and no flux turned negative. Because the deep store can
be drawn below zero, its drainage term, and with it total runoff, turns slightly negative on a few days (minimum
$-2.5$ to $-10$\,mm per day in the cells tested). The deep store's drainage was added after a first version let it accumulate 2--3\,m of water
over 81 years, ten times the original's change.

\paragraph{Training.} The networks are trained one step ahead (yesterday's stores from the original) on 40 million
sampled cell-days of 730 training cells, with 80 further cells for early stopping, minimising the sum over the eight
targets of the squared errors standardised by each target's spread (Adam with a one-cycle learning-rate schedule peaking at $3\times10^{-3}$, batches of
16\,384, ten epochs; about 15 minutes per member on one GPU). They are then run freely: the stores are carried by the
model over 81 years. A free run whose error grew to twice the one-step error would have triggered rollout training; it
never did. Specific humidity was dropped as an input because the original's daily exchange does not read it; removing
it changed no score beyond run-to-run noise.

\subsection{Yearly climate statistics}

At the end of each year, LandEmu computes 25 statistics of the year's daily weather: mean, coldest-month and
warmest-month temperature; degree days above 5\,$^\circ$C; frost days and days above 30\,$^\circ$C; annual and
summer precipitation; potential evaporation; the annual, summer and driest three-month climatic water balance; the
longest dry spell; mean relative humidity; summer vapour pressure deficit and its 10-day extremes; mean shortwave
and net longwave radiation; and, for each of the seven tree types, the number of days outside its temperature-stress
range. These enter a trailing 30-year window, whose means are the climate the forest maps see.

\subsection{Yearly forest maps}\label{sec:maps}

For each yearly quantity $q$ (six structure quantities and twelve stand quantities) a separate gradient-boosted tree
model \citep{Ke2017} predicts
\begin{equation}
  q_{c,y} = A^{q}_{c,y} + g_q\!\left(\bar{x}_{c,y},\; \bar{x}_{c,y} - \bar{x}^{A}_{c,y},\; A_{c,y},\;
  s^{0}_{c},\; \Delta s^{0}_{c},\; \phi_c,\lambda_c,\sigma_c\right),
\end{equation}
where $A^{q}_{c,y}$ is the \emph{anchor}: the original's mean of $q$ at cell $c$ in year $y$ over its training runs
and climates; $\bar{x}_{c,y}$ is the trailing 30-year climate and $\bar{x}^{A}_{c,y}$ the corresponding climate of the
anchor runs, so that their difference measures how far this climate departs from what the anchor represents;
$A_{c,y}$ collects all anchors of the cell; $s^{0}_{c}$ is the run's own 2000--2019 state and $\Delta s^{0}_{c}$ its
offset from the training runs' mean (the known start state); and $\phi$, $\lambda$, $\sigma$ are latitude, longitude
and soil code. Biomass per tree is predicted as a logarithmic ratio to its anchor, which weights relative errors
equally across dense and sparse stands; the other quantities as differences. Each model has 400 trees with 31 leaves,
learning rate 0.05 and 80\,\% row and column subsampling, and is trained on every fourth year of the training runs:
five runs under nine climate models in the test folds (about three million cell-years), all six runs under all ten
climate models in the deployment release (about four million).

Two properties of this form matter. First, the map is not recursive: no prediction is fed back into the forest
maps, so a year's error cannot carry into the next (P3). The only quantity passed from the maps to the daily model is
the stand, which affects next year's fluxes but not next year's maps. Second, the anchor carries what is specific to a
cell --- its soil, its history, the composition the original grew there --- so the learned part only has to explain how
a cell departs from it under a different climate. The price is that the anchor must exist: the original must have
been run on every cell the emulator is to run on (Sect.~\ref{sec:limits}). For the full-grid application
(Sect.~\ref{sec:global}) the anchor is instead the cell's own simulated 1985--2014 state, constant in time.

\subsection{Implementation}\label{sec:impl}

LandEmu is a Python package whose only run-time dependencies are numpy and LightGBM; the networks are ported from
PyTorch to numpy (maximum difference $10^{-6}$ of each output's range on 20\,000 real input rows). The yearly climate
statistics of the package equal those of the training pipeline exactly, and a validation release reproduces the
scored daily fluxes of the evaluation to rounding (median per-cell difference of 30-year means 1.6--2.9$\times10^{-4}$)
and the yearly quantities with a median difference of zero (99th percentile 0.1--7\,\%). The remaining yearly
differences have one cause: the training table converted the forcing, stored in 0.1\,$^\circ$C steps, in single
precision, so days at exactly 30.0\,$^\circ$C counted as hot days; continuous forcing from a host model never sits
exactly on such a threshold. A release (models plus per-cell set-ups for 1\,050 cells) occupies 183\,MB. The interface
offers a whole year (\texttt{step\_year}) or a day at a time (\texttt{begin\_year}, 365\,\x\,\texttt{step\_day},
\texttt{end\_year}, bit-identical to the yearly call), so that a host model can exchange fluxes daily.

\subsection{Variant with a \texorpdfstring{CO$_2$}{CO2} input}\label{sec:co2var}

For LPJmL-FIT with nitrogen limitation and rising CO$_2$, the same architecture receives the year's CO$_2$
concentration as an additional daily-network input, and the forest maps receive the trailing 30-year mean CO$_2$ and
its change since 1990--2019. It is trained on runs of that model (Sect.~\ref{sec:resp}).

\section{Training data and evaluation}\label{sec:eval}

\subsection{Simulations of the original}

The main data set is a panel of 1\,050 cells: 105 blocks of ten contiguous 0.5$^\circ$ cells, chosen to balance 20
climate strata (Fig.~\ref{fig:design}a). Each block was spun up for 1000 years and run over 2000--2019 under observed
weather (GSWP3-W5E5; \citealt{Lange2021W5E5}) and over 2020--2100 under bias-adjusted output \citep{Lange2019} of ten
CMIP6 climate models (GFDL-ESM4, IPSL-CM6A-LR, MPI-ESM1-2-HR, MRI-ESM2-0, UKESM1-0-LL, CanESM5, CNRM-CM6-1, CNRM-ESM2-1,
EC-Earth3, MIROC6) for three scenarios (SSP1-2.6, SSP3-7.0, SSP5-8.5; \citealt{ONeill2016}), plus a control with
observed 1990--2019 years resampled. We call each combination of climate model and scenario a \emph{leg}. Six
independent runs exist, each with its own spin-up and random seed, all with 25 patches, nitrogen limitation off and
CO$_2$ constant after 2019. Per-tree output is written every year (trees taller than 5\,m, as in the original's
standard output); daily fluxes are written for 13 legs of five climate models and for the control, in runs 1--4.

\subsection{Held-out tests}

The forest maps are tested in folds that hold out one of five climate models entirely, together with run 4: they are
trained on the other nine climate models and runs 1, 2, 3, 5 and 6, and run 4 under the held-out model is the truth
(Fig.~\ref{fig:design}b). With three scenarios this gives 15 test cases. The cells themselves are seen in training,
as they would be in deployment; this choice is discussed in Sect.~\ref{sec:limits}. The daily networks are trained on
runs 1--4 under all 14 legs with daily output and tested on 24 blocks (240 cells, 238 with GPP of at least
50\,gC\,m$^{-2}$\,yr$^{-1}$) never used in their training, including the five biome blocks of Fig.~\ref{fig:runtime}.
For them, therefore, only the place is new: they have seen every climate model they are scored on. In the full
emulator the daily networks read the stand predicted by forest maps that never saw the test run or the test climate
model. Earlier tests with a climate model withheld from the daily networks are summarised in Sect.~\ref{sec:fluxes}.

\subsection{Measuring error against the original's own variability}

Because a second run of LPJmL-FIT differs from the first, a fixed tolerance such as ``within 10\,\%'' is uninformative:
two runs of the original agree on all six structure quantities within 10\,\% (or within the typical two-run spread of
the cell's density class, where that is larger) in only 17--21\,\% of cells. Following
practice in climate-model emulation and model verification
\citep{WatsonParris2022,Baker2015,Lutjens2025}, we instead compare the distribution of errors over cells with that of
a further run. For cell $c$ and quantity $q$ against the unseen truth run $T$, the emulator's error is
$e_E(c)=|E(c)-T(c)|/|T(c)|$ and a further run $R$'s error is $e_R(c)=|R(c)-T(c)|/|T(c)|$, pooled over the available
further runs. The \emph{error ratio} at centile $k$ is
\begin{equation}
  \rho_k = \frac{P_k\left[e_E\right]}{P_k\left[e_R\right]},
\end{equation}
with $k=50$ (the typical cell) and $k=90$ (the worst tenth of cells). A further run of the original scores
$\rho\approx1$ by construction; the mean of three runs scores about 0.82, and a perfect predictor of the expectation
about 0.71, because an average beats a single random realisation. Values below one are therefore possible and, for a model that predicts the expectation, expected. We
apply the ratio to 30-year means (2071--2100 for structure; 2021--2050 and 2071--2100 for fluxes), to the climate
response (late minus early period), and to single years (for drift). Area totals (cosine-latitude weighted) are
reported as signed relative errors, because a coherent bias of a few per cent passes every per-cell test yet does not
average out in a sum.

\paragraph{Pass standard.} Forests with at least five trees per patch (``dense''; about 80\,\% of tree-bearing cells)
must have a typical error ratio of at most 1.2, a worst-tenth ratio of at most 1.5 and area totals within 5\,\%;
sparse stands must stay within 3 times a further run and their totals within 25\,\%. A criterion passes when its
median over cases meets the limit and at least 12 of 15 cases do. The standard as first written classed a cell as dense
by the held-out run itself. That selects cells on the outcome being predicted --- a cell whose truth run happens to
lose trees moves into the sparse class --- so we class cells by the anchor instead, before the run, and report the
original definition beside it (Sect.~\ref{sec:structure}). The daily fluxes need not be as close to the original as a
second run is; for them we propose a typical and worst-tenth ratio of at most 3 (the sparse-stand limit above), a
climate-response ratio of at most 1.5 and area totals within 5\,\%, numbers set after the first daily scores were
known. \emph{Drift} means the gap to the original growing, not the values changing; we measure it as the ratio of the
single-year typical error ratio in 2081--2100 to that in 2021--2040, with a limit of 1.2. For the forest structure,
whose early years benefit from the known start state, the pre-registered form of this rule failed for that reason, and
we additionally report the ratio to 2051--2070, a baseline chosen after that result (Sect.~\ref{sec:drift}). Because
several of these choices were made with some scores known, we report every value so that readers can apply their own
limits.

\paragraph{Reference arms.} Every score is reported beside (i) a further run scored as if it were the emulator,
which must give about 1 and checks the scorer; (ii) the mean of other runs, a reference for a near-ideal predictor of
the expectation; (iii) the anchor alone, without the learned departure; and (iv) the run's 2019 state carried
forward, which shows whether the measure can detect drift.

\begin{figure}[!tbp]
  \centering
  \includegraphics[width=\textwidth]{\figpath/fig_design.pdf}
  \caption{Training and test design. (a) The 1\,050 panel cells (105 blocks of ten cells), the 24 blocks never used to
  train the daily networks, and the five biome blocks. (b) One of the five test folds of the yearly forest maps: one
  climate model (here GFDL-ESM4) and one independent run are held out; further runs under the held-out climate model
  give the yardstick, and one of them is scored as if it were the emulator to check the scorer.}
  \label{fig:design}
\end{figure}

\section{Results}\label{sec:results}

\subsection{How the forest maps behave at a single cell}

Figure~\ref{fig:cell} shows what the yearly maps do. At the temperate and the boreal cell the original's runs scatter
widely: its unseen run and a further run of the same climate wander apart by up to two trees per patch and by 50\,\%
in biomass per tree for decades, because a few deaths or recruitments in 25 patches change the stand. The anchor ---
the mean of the training runs under other climates --- smooths this noise but knows nothing about the held-out
climate. The emulator follows the anchor where the climate resembles the training climates and departs from it where
it does not: at the temperate cell it keeps biomass per tree above the declining anchor in the second half of the
century, as the unseen run does. It does not reproduce the excursions of the single run, and it should not: they are
not predictable from climate.

\begin{figure}[!tbp]
  \centering
  \includegraphics[width=\textwidth]{\figpath/fig_cell_example.pdf}
  \caption{The yearly forest maps at two single cells under a held-out climate model (GFDL-ESM4, SSP3-7.0): (a, b)
  trees taller than 5\,m per patch and (c, d) above-ground biomass per tree. Grey: training runs under the nine other
  climate models; violet: their mean, the anchor; black: the unseen run of LPJmL-FIT; orange: a further run under the
  same climate; blue: the emulator.}
  \label{fig:cell}
\end{figure}

\subsection{Forest structure}\label{sec:structure}

Over all 15 held-out cases (Fig.~\ref{fig:structure}), the dense forests' typical error of tree density is
1.16 times that of a further run (limit 1.2; met in 11 of 15 cases, one short of the required 12), and that of biomass
per tree 1.07 times (14 of 15). The four trait medians are closer to the unseen run than a further run is
(0.60--0.83 times): the emulator predicts the expectation, while a further run is itself a random draw. In the worst
tenth of dense cells tree density reaches 1.29 times and biomass per tree 1.01 times a further run (limit 1.5). Area
totals over dense cells are off by a median of $-0.1$\,\% for stems (cases $-2.0$ to $+1.2$\,\%) and $-1.5$\,\% for
biomass ($-5.0$ to $+3.4$\,\%; one case just past the 5\,\% limit). Sparse stands are harder: tree density 2.74 times a
further run (limit 3), biomass per tree 1.34 times, stem and biomass totals a median $-2.3$\,\% ($-13$ to $+12$\,\%)
and $-5.7$\,\% ($-18.5$ to $+6.6$\,\%; limit 25\,\%). Of the 22 criteria of the pass standard, 21 are met. With dense
cells classed by the held-out run, as the standard was first written, 19 of 22 are met: tree density then fails in
typical dense cells (1.12, 9 of 15 cases) and typical sparse cells (2.95, 8 of 15), and the sparse stem total is
$+17.7$\,\% (11 of 15).

\begin{figure}[!tbp]
  \centering
  \includegraphics[width=\textwidth]{\figpath/fig05_structure.pdf}
  \caption{Forest structure in 2071--2100 under a held-out climate model and run (15 cases: five climate models
  $\times$ three scenarios). (a--c) Error ratio for six quantities: typical dense cell, worst tenth of dense cells and
  typical sparse cell; dots are cases, bars medians. Orange: a further run of LPJmL-FIT scored the same way, which must
  give about 1. (d) Signed area-total errors of stems and biomass. Dashed: limits of the pass standard.}
  \label{fig:structure}
\end{figure}

\subsection{No drift in the forest state}\label{sec:drift}

Run year by year from 2020 to 2100 (Fig.~\ref{fig:drift}), the emulator's error starts below that of a further run,
because it knows the run's own 2000--2019 state, and rises over about three decades as that knowledge fades. It then
levels off: biomass per tree at about one further run's error from the 2050s on, tree density at 1.10--1.12 in the
2050s and 1.13--1.21 in the 2090s. From 2051--2070 to 2081--2100 --- the baseline chosen after the pre-registered one
failed, see below --- the typical error grows by a factor of 1.05 for both quantities (all 15 cases below the limit of
1.2); the trait medians grow by 1.10--1.19 while remaining at 0.62--0.85
times a further run. The rise in the first decades does not depend on the scenario (tree density 1.17 under SSP1-2.6,
1.20 under SSP5-8.5), which identifies it as the fading start state rather than climate extrapolation. What remains is
a slow creep of tree density in the second half of the century that increases with warming (+0.01, +0.07 and +0.12
from the 2050s to the 2090s for SSP1-2.6, SSP3-7.0 and SSP5-8.5). The reference arms behave as they should: a further
run stays at 1.0, the mean of other runs at 0.8, the anchor alone rises steadily to 1.3--1.6, and the carried-forward
2019 state rises to more than twice a further run's error, so the measure does detect drift. The pre-registered drift
rule, which compares 2081--2100 with 2021--2040, fails (1.17 for tree density, 1.32 for biomass per tree; 8 and 2 of 15
cases within the limit). We attribute this to the artificially low early error, which the scenario-independence of the
early rise supports, but the rule as written is failed, and the late creep of tree density under strong warming is a
real, if slow, growth of the gap.

\begin{figure}[!tbp]
  \centering
  \includegraphics[width=\textwidth]{\figpath/fig07_drift.pdf}
  \caption{The forest maps run year by year: typical-cell error ratio of dense cells per year, median over the 15
  held-out cases (shading: 10th--90th percentile of cases for the emulator), for (a) tree density, (b) biomass per tree
  and (c) median wood density. Reference arms: a further run (orange), the mean of five other runs (green), the anchor
  without the learned departure (violet) and the 2019 state carried forward (grey dashed).}
  \label{fig:drift}
\end{figure}

\subsection{Daily carbon and water exchange}\label{sec:fluxes}

At the three biome cells of Fig.~\ref{fig:daily} the full emulator --- daily networks never trained on these cells,
fed a stand from forest maps that never saw the climate model or the run --- reproduces the seasonal course and the
day-to-day variability of GPP and evapotranspiration: the spring onset in the temperate and boreal forests, the short
wet season at the semi-arid cell, the response to individual dry spells. Year-to-year variations of annual GPP follow
the original's at the temperate and semi-arid cells (correlation 0.93 and 0.94; a further run: 0.91 and 0.96), less so
at the boreal cell (0.66; a further run: 0.99). The 2021--2100 means differ by +0.8\,\% at the temperate and
$-0.6$\,\% at the boreal cell, but by +12\,\% at the semi-arid cell, where the original's own two runs differ by
3\,\%. The top-metre water store is the
least faithful daily variable: at the temperate cell it dries 20\,mm too deep in summer, and at the boreal cell it
holds up to 50\,mm more water in summer than the original, while the fluxes stay close. This is expected: the
networks are trained on fluxes and store changes, and the store is a state integrated by the bucket, not a
prediction target.

Over the 238 held-out cells and 13 legs (Fig.~\ref{fig:dailyscores}), the typical error of 30-year mean
GPP is 2.4 times a further run's (median over legs; limit 3, met on 13 of 13 legs), evapotranspiration 3.1 times
(5 of 13), NPP 3.9 times (0 of 13) and runoff 5.2 times. Since two runs of the original differ by only 1.0--1.6\,\% in
a cell's 30-year GPP, evapotranspiration or NPP, this means typical per-cell errors of about 3--6\,\% for GPP and
evapotranspiration and 5--6\,\% for NPP. Area totals are off by a median of 1\,\% for GPP, NPP and evapotranspiration
and 2\,\% for runoff, and by at most 2.5\,\% on any leg. The climate response (2071--2100 minus 2021--2050) is 1.8--2.5 times a further run's error, above
the proposed 1.5. Runoff is hardest because it is a small residual of precipitation minus evapotranspiration, so the
evapotranspiration error passes into it amplified; NPP is hard because it is a small difference of two large terms in
the most productive cells.

The emulated stand costs accuracy. Fed the original's own stand instead, the same daily networks reach 2.2 (GPP), 2.6
(evapotranspiration), 3.4 (NPP) and 4.5 (runoff) times a further run's error. Compared leg by leg and period by
period, the emulated stand adds a median 13\,\% to the GPP error and 35\,\% to the evapotranspiration error, most of it
through the cover of the plant types. On these legs no flux drifts away from the original: the late-to-early ratio of
single-year errors is 0.90--1.03 for the medians of all four fluxes, and only GPP on the two warmest legs
(UKESM1-0-LL SSP3-7.0, IPSL-CM6A-LR SSP5-8.5) exceeds the limit (1.26).
The deep water store stays within the range of the original's change, and over 81 years the across-cell mean of the
top-metre store stays between 80 and 91\,mm.

In this evaluation the climate models are seen by the daily networks; only the cells are new. That matters: in
earlier tests with a climate model withheld from daily-network training entirely, the gap to the original in GPP and
NPP grew by a factor of 1.3--1.4 over the century. The cause was the daily carbon response degrading as the climate
moved away from the training climates, not the water stores, since a run with the stores reset daily from the original
drifted identically. Adding further climate models to training did not reduce this growth beyond run-to-run noise for
UKESM1-0-LL; it falls (to 1.17 for NPP) only when that model itself is in training, and even then the full emulator's
GPP gap grows by 1.26 on that model's warmest leg. The daily networks therefore have not been shown to hold their
accuracy under a climate model they have not seen.

\begin{figure}[!tbp]
  \centering
  \includegraphics[width=\textwidth]{\figpath/fig_daily_series.pdf}
  \caption{Daily output of the full emulator at three biome cells never used to train the daily networks
  (GFDL-ESM4, SSP3-7.0; forest maps with this climate model and run held out). (a--c) Daily GPP, (d--f)
  evapotranspiration and (g--i) top-metre soil water in 2085; (j--l) annual GPP 2021--2100. Black: the unseen run of
  LPJmL-FIT; orange: a further run; blue: the emulator.}
  \label{fig:daily}
\end{figure}

\begin{figure}[!tbp]
  \centering
  \includegraphics[width=\textwidth]{\figpath/fig_daily_scores.pdf}
  \caption{Daily fluxes over 238 held-out cells and 13 climate-model legs (dots: legs; bars: medians). (a) Typical-cell
  error ratio of 30-year means (larger of 2021--2050 and 2071--2100); (b) error ratio of the climate response; (c)
  drift, the late-to-early ratio of single-year errors (1 = no change in the gap to the original); (d) area-total
  errors. Blue: the full emulator; green: the same daily networks fed the original's stand. Orange line: a further run
  of the original; dashed: proposed limits.}
  \label{fig:dailyscores}
\end{figure}

\subsection{Learning the warming response}\label{sec:resp}

The warming response is the property an ESM needs most and the one most emulators never test. We tested it on
LPJmL-FIT with nitrogen limitation and rising CO$_2$, run on 674 cells that each represent one region of a 674-region
climate classification of the land surface (25 patches each), under the ten climate models and three scenarios, with
CO$_2$-fixed twins of SSP3-7.0 for the training climate models. Two climate models of low and high climate sensitivity,
MPI-ESM1-2-HR and IPSL-CM6A-LR, were held out entirely, together with 20\,\% of the cells; both lie within the warming
range of the eight training models, of which CanESM5 and UKESM1-0-LL warm more, so the test climate is covered by the
training climates. An independent second run of the
original exists for the held-out legs. Here the forest maps are anchored on each cell's own 1990--2019 state, so no
future output of the original enters the prediction. The pass criteria were registered before the runs: on every
held-out leg a global change of the same sign as and within $\pm25$\,\% of the original's, a per-cell regression slope between 0.8 and 1.2,
and a per-cell root-mean-square error at most 1.5 times the second run's.

With rising CO$_2$, the original's vegetation carbon increases by 32--45\,\% by 2071--2100 on the held-out legs
(Fig.~\ref{fig:resp}a), mostly through CO$_2$ fertilisation. The emulator reproduces 81--97\,\% of this change, while
the original's independent second run reproduces 61--89\,\%: the original's own global response varies by up to
40\,\% between runs on this grid, because single cells stand for whole regions. Trained only on SSP1-2.6, the same maps
reproduce 37--43\,\% --- the failure seen when an emulator meets a climate warmer than its training climates.

With CO$_2$ fixed at its 2019 level, the original loses vegetation carbon under warming, by up to 20.7\,\% for
IPSL-CM6A-LR under SSP5-8.5 (second run: 26.3\,\%). The maps as registered reproduce 67\,\% of the IPSL-CM6A-LR loss
under SSP3-7.0 and 57\,\% under SSP5-8.5. After seeing this result we added CO$_2$-fixed twins of SSP1-2.6 and SSP5-8.5
to training, so that warming and CO$_2$ rise could be told apart for strong warming; the held-out models and cells were
unchanged and nothing else was tuned. This raised both fractions to 72\,\%, short of the 75\,\% we had stated as the
expectation for this step, and left the rising-CO$_2$ results unchanged within 0.01. Figure~\ref{fig:resp} shows this
second version. The SSP1-2.6-only maps predict gains on all four warming-only legs, because all their training legs
carry rising CO$_2$. Per cell (Fig.~\ref{fig:resp}b) the emulator's change is as close to the original's as the second
run's is (root-mean-square error 0.65--0.80 times the second run's), and the slopes of both against the original are
below one (0.73 and 0.75): the original's per-cell change contains run-to-run noise, which flattens the regression slope
of any predictor, including a perfect estimate of the expected response.

For GPP the picture splits. The daily networks fed the original's stand reproduce the global GPP change within 4\,\%
and the per-cell change with slopes of 0.92--0.99, with rising and with fixed CO$_2$. Fed the emulated stand, the
global change with rising CO$_2$ is still within 6\,\%, but the per-cell slope falls to 0.66--0.79, because the forest
maps' response of leaf area and tree cover is too uniform across cells (slopes 0.73--0.83, against 0.89--0.96 for the
second run). The global totals are right; the cell-to-cell spread of the response is too small. On the held-out cells,
a learning curve over 135, 269 and 539 training cells reduced this error by 6--12\,\% per doubling, which points to the
number of training cells rather than the model form as the limit; extrapolated, reaching 1.2 times a second run for leaf
area would take about 12\,000 training cells, the scale of a full-grid training set.

Judged by the registered criteria, the test was \emph{not passed}. The global changes and the root-mean-square errors
pass on every leg. The slope criterion fails for vegetation carbon on three of four rising-CO$_2$ legs --- a criterion
that the original's own second run also fails (slopes 0.72--0.84), for the reason given above --- and for GPP fed the
emulated stand on all four, which is a real defect. With CO$_2$ fixed, the largest warming-only loss stays damped by
about 30\,\%.

\begin{figure}[!tbp]
  \centering
  \includegraphics[width=\textwidth]{\figpath/fig_response.pdf}
  \caption{The warming response under two climate models never seen in training (LPJmL-FIT with nitrogen limitation,
  674 representative cells). (a) Area-weighted change of vegetation carbon from 1990--2019 to 2071--2100 with rising
  CO$_2$ and with CO$_2$ fixed at 2019: the original, its independent second run, the emulator (the version with
  CO$_2$-fixed twins of every training scenario, added after the first held-out results; see text), and the same
  emulator trained on SSP1-2.6 only. (b) Per-cell change for IPSL-CM6A-LR, SSP5-8.5, fixed CO$_2$: emulator (blue) and second
  run (orange) against the original; dashed lines are least-squares slopes, the solid line is the 1:1 line.}
  \label{fig:resp}
\end{figure}

\subsection{Full-grid application}\label{sec:global}

To show that the emulator runs globally, we retrained the forest maps on four runs of LPJmL-FIT (nitrogen off) under
GFDL-ESM4 --- historical, SSP1-2.6, SSP2-4.5, SSP3-7.0 --- made with a different build of the model, anchored each cell
on its own simulated 1985--2014 state, and emulated each future scenario with maps that never saw it. The daily
networks were not retrained (these runs store no daily output). LandEmu then ran on all 64\,240 land cells from 1975 to
2100 (Fig.~\ref{fig:global}). Over 1985--2014 the emulated global GPP equals the original's (+0.04\,\%), NPP is
1.1\,\% high, transpiration 8.5\,\% low and runoff 5.8\,\% high. Benchmarked against observations with the ILAMB
system, the emulator's overall scores are within 0.03 of the original's on all 27 data-set comparisons; for biomass,
leaf area and cover this is not a test of the emulator, because its historical forest state is anchored on the
original's. Over the same
period the original's GPP rises from 151 to 162\,PgC\,yr$^{-1}$ because CO$_2$ rises in its historical forcing, while the
emulator's stays flat: the nitrogen-off emulator has no CO$_2$ input by design (Sect.~\ref{sec:target}), so it
reproduces the level but not this trend.

Under the unseen scenarios, the typical per-cell error in 2071--2100 is 6.0--6.4\,\% for GPP, 4.4--4.8\,\% for
evapotranspiration, 11--15\,\% for vegetation carbon and 7--8\,\% for dense-cell tree density, always smaller than when
the 1985--2014 state is carried forward (11--15\,\%, 7--12\,\%, 18--26\,\% and 10--11\,\%). The global responses agree
in sign in 13 of 15 cases (five quantities, three scenarios), but the carbon responses are 3--5 percentage points too
low, and under SSP3-7.0, the warmest scenario and the only one outside the training range, the emulator captures only
a third of the original's vegetation-carbon loss ($-2.8$ against $-9.2$\,\%). This is the extrapolation failure of
Sect.~\ref{sec:resp} reproduced on the full grid: with only milder scenarios in training, the strongest response is
damped. On an earlier set of eight independent runs of this model under GFDL-ESM4, scored over all 57\,726 tree-bearing cells against a further run, an earlier form of the forest maps that
predicts 30-year windows rather than single years passes the dense-cell standard when only the run is held out (tree
density and biomass per tree 0.92 times a further run) but not when the scenario is held out as well (1.24 and 1.47
times).

\begin{figure}[!tbp]
  \centering
  \includegraphics[width=\textwidth]{\figpath/fig_global.pdf}
  \caption{LandEmu on the full 0.5$^\circ$ grid. Vegetation carbon in 2071--2100 under SSP3-7.0 in (a) LPJmL-FIT and
  (b) the emulator, whose forest maps never saw SSP3-7.0, and (c) their difference (cells with at least
  0.5\,kgC\,m$^{-2}$). (d) Global vegetation carbon and (e) global GPP as change relative to the original's 1985--2014
  mean (5-year running means): LPJmL-FIT solid, emulator dashed. The historical forest state is anchored and therefore
  not a test.}
  \label{fig:global}
\end{figure}

\subsection{Cost}\label{sec:cost}

Measured end to end on one core over 1\,050 cells and 81 years, the whole emulator costs
$7.6\times10^{-3}$ core-seconds per cell-year (Fig.~\ref{fig:speed}), 97\,\% of it in the daily networks. LPJmL-FIT,
timed on one core of the same computing cluster at the temperate cell (marginal rate), costs 0.267 core-seconds at 25
patches and 7.06 at 500 patches, so the emulator needs 35 times fewer core-seconds than the original at its training
setting and about 900 times fewer than at the publication setting. The original's cost depends on how many trees a
cell carries: across the five biome cells it is 0.20--0.33 core-seconds at 25 patches, which gives factors of 26--44. The re-implemented differentiable physics of the first design cost
1.23 core-seconds. The emulator lies inside the allowance for an intermediate-resolution atmosphere by a factor of 1.8
(T31) to 4 (T63). The nitrogen-on variant costs $7.3\times10^{-3}$ core-seconds, 115--130 times less than that model
on the same cells. Training is cheap as well: about 15 GPU-minutes per network and seconds to minutes per forest map.
The cost is dominated by one matrix product per day and network, so batching over years or a GPU would reduce it
further; neither is used here.

\begin{figure}[!tbp]
  \centering
  \includegraphics[width=0.66\textwidth]{\figpath/fig_speed.pdf}
  \caption{Cost per cell-year on one core (marginal rate). LPJmL-FIT at its training and publication settings, the
  differentiable re-implementation of its daily physics, and LandEmu (dark blue: daily networks, 97\,\%; light blue:
  yearly statistics and forest maps). Dashed: land allowances for T31 and T63 atmospheres (10\,\% of a coupled atmosphere's cost).}
  \label{fig:speed}
\end{figure}

\section{Discussion}\label{sec:disc}

\subsection{LandEmu against the requirements for ESM-ready land emulators}

Table~\ref{tab:req} compares LandEmu with the six requirements derived in the review \citep{PriesnerReview}. Three
design choices carry most of the result. \emph{Learning the whole daily step} is what makes the emulator fast: the
literature's large speed-ups come from replacing an expensive iterative solver or a whole model step
\citep{R00067,R02497,R00568}, and our profile of LPJmL-FIT shows that nothing short of this can give an order of
magnitude. Re-implementing physics can be faster than legacy code when the code is a column model without a
per-individual loop \citep{R00975}; for an individual-tree model it was not. \emph{Predicting the expectation} removes
the patch ensemble, which is a variance-reduction device of the original, not a process; at the publication setting it
is worth more than every process combined. \emph{Keeping the forest state non-recursive} sidesteps the error
accumulation that made our one-step demography model unusable and that rollout losses \citep{R00593,R01133} or
trajectory fitting \citep{R00661} address in other emulators. The daily stores are recursive, but they are closed
physical stores rather than a learned latent state, and they did not walk away from the original in 81-year free runs.

\begin{table}[!tbp]
\centering
\caption{LandEmu against the six requirements for ESM-ready land emulators of \citet{PriesnerReview}.}\label{tab:req}
\small
\begin{tabularx}{\textwidth}{@{}p{3.3cm}X@{}}
\toprule
\textbf{Requirement} & \textbf{LandEmu v1} \\
\midrule
Prognostic state and memory & explicit snow, top-metre and deep water stores; trailing weather windows; 30-year
climate window; stand carried yearly. No learned latent state. Forest state not recursive by design. \\
Coupling interface & daily latent-heat-relevant fluxes (interception, transpiration, soil evaporation), runoff, snow;
GPP and NPP; leaf area and plant-type cover from which albedo and roughness can be derived. Missing: energy fluxes and
surface temperature (LPJmL-FIT computes none), heterotrophic respiration and fire, hence the net CO$_2$ flux. \\
Conservation & water closed exactly every day by construction. Carbon is not tracked between NPP and vegetation carbon,
which are predicted separately. \\
Stability under feedback & structure error grows by a factor of 1.05 after the start state fades, with a slow creep of
tree density under strong warming (the pre-registered drift rule fails); no flux drift on climate models seen in
training except GPP on the two warmest legs; day-by-day interface identical to the yearly call. No feedback with a
host model tested here. \\
Out-of-distribution climate & yearly maps: climate model and independent run held out; daily networks: places held out,
climate models seen (with one withheld, the GPP and NPP gap grew 1.3--1.4-fold); warming response tested on unseen
climate models with and without CO$_2$ rise, not passed as pre-registered; extrapolation beyond the training scenarios
fails (one third of the loss). \\
Uncertainty, interfaces, openness & errors expressed relative to the original's run-to-run variability; deterministic
prediction of the expectation, no predictive distribution; numpy + LightGBM package with daily and yearly calls. \\
\bottomrule
\end{tabularx}
\end{table}

\subsection{What limits the fidelity}

Three findings locate the remaining error. First, the emulator is \emph{data-limited in the space of climates, not in
the number of runs}. On the panel, adding five climate models to training raised the share of cells close to the
unseen run by 0.023 and lowered the biomass-per-tree error ratio from 1.64 to 1.42, while adding two runs of the
original raised it by only 0.005. The warming-response test makes the same point from the other side: with training
climates that cover the test climate the response is reproduced, with only milder climates it is damped by more than
half. Second, \emph{sparse stands} carry most of the structure error. Reweighting the loss towards them (logarithmic
targets, Poisson and Tweedie objectives) only moved error between sparse and dense cells, while the mean of three runs
is 0.82 times a further run in every density class --- so their response is predictable in principle, and what is
missing is information about how sparse cells respond. Third, the daily networks' own error, two to three times a
further run, still exceeds the 10--40\,\% that the emulated stand adds, but the stand's share is concentrated in
evapotranspiration and in the cell-to-cell spread of the climate response.

\subsection{Limitations}\label{sec:limits}

\begin{itemize}
  \item \textbf{The original must have been run on every emulated cell.} The forest maps depart from a per-cell
        anchor. In the panel version the anchor is the original's mean over its training runs for each year to 2100;
        in the global version it is the cell's simulated 1985--2014 state. Either way, a cell never simulated by
        LPJmL-FIT cannot be emulated, and on the response test the error on the 20\,\% of held-out cells is 1.8--2.2
        times a second run's for leaf area and cover, compared with near-parity on training cells.
  \item \textbf{The cells are seen in training.} The structure tests hold out the climate model and the run, not the
        place. This mirrors deployment on the cells the original was run on; with cells held out as well, the panel
        scores were clearly worse.
  \item \textbf{Daily NPP and runoff} miss the proposed limit (3.9 and 5.2 times a further run), and the climate
        response of every flux is 1.8--2.5 times a further run's error.
  \item \textbf{The daily networks have seen every climate model they were scored on.} With a climate model withheld,
        their GPP and NPP gap to the original grew 1.3--1.4-fold over the century.
  \item \textbf{Some evaluation choices were made with results known}: classing dense cells by the anchor, the
        2051--2070 drift baseline for the forest structure, the limits for the daily fluxes, and the CO$_2$-fixed
        training runs of the response test. We report the outcome under the pre-registered choices beside each.
  \item \textbf{Extrapolation.} Under a scenario warmer than all training scenarios the response is damped
        (Sect.~\ref{sec:global}); deployment should keep the training climates broader than the target climate.
  \item \textbf{Scope of the outputs.} The structure targets cover trees taller than 5\,m, as the original's standard
        output does, and the traits are represented by their medians, not their full distributions. There are no
        energy fluxes and no net CO$_2$ flux.
  \item \textbf{Scale of the evaluation.} The detailed evaluation uses 1\,050 of 54\,020 tree-bearing cells; the
        full-grid application has one run per scenario, so its errors cannot be set against a further run.
\end{itemize}

\subsection{Outlook}

The day-by-day interface makes LandEmu callable from an atmosphere model; experiments with Earth system models will be
reported separately. The most direct improvements are training on a full-grid set of runs that covers warmer
climates (which removes the unseen-cell limit and the extrapolation problem together), a stand map aimed directly at
the daily networks' inputs, explicit litter and soil carbon pools so that the net CO$_2$ flux is complete, and a
calibrated predictive spread scored against several independent runs of the original.

\section{Conclusions}\label{sec:concl}

LandEmu v1 emulates an individual-tree, trait-based DGVM as a whole: its daily carbon and water exchange and its yearly
forest structure, trait composition and stand, from daily weather alone, with water conserved exactly. Its design ---
learn the whole daily step, predict the expectation over the patch ensemble, keep the forest state non-recursive,
train on climate contrast and judge against the original's own variability --- followed from measurements that ruled
out a hybrid of learned demography and re-implemented physics. Under a held-out climate model and run, dense forests
are reproduced close to the original's run-to-run variability (tree density 1.16, biomass per tree 1.07 times a further
run; trait medians better than a further run), and the structure error grows only slowly once the start state has
faded. On places held out from their training, the daily networks reproduce GPP and evapotranspiration to a few per
cent per cell, without drift on the climate models they were trained on; under an unseen climate model their carbon
fluxes drift. Trained on climates that cover the test climate, the emulator reproduces the change in vegetation carbon
under rising CO$_2$ as well as a second run of the original does on the global total, but it damps the largest
warming-only losses, the cell-to-cell spread of the response and the response of sparse stands, and it did not pass
the test as pre-registered. At $7.6\times10^{-3}$ core-seconds per cell-year it costs 35--900 times less than the
original and fits the land budget of an intermediate-resolution Earth system model.

\section*{Code and data availability}
The emulator package (\texttt{python/src/lpjmlfit\_emulator/landemu}), the training, scoring and release scripts, the
scripts that produce every figure of this paper (\texttt{scripts/paper}) and the decision records behind every number
are in the project repository \url{https://github.com/rimajj/LPJmLFIT_Emulator} (to be archived on Zenodo on
acceptance). The LPJmL-FIT simulations and the released emulator (models and per-cell set-ups) are stored at the
Potsdam Institute for Climate Impact Research and will be archived with the code.

\section*{Author contributions}
JP conceived the study, defined its goals and acceptance standards, and made the design decisions. The implementation,
simulations, analysis, figures and drafting were carried out by an AI research assistant (Claude, Anthropic) under JP's
direction, with every decision recorded in the repository; JP is responsible for the final content.

\section*{Competing interests}
The author declares no competing interests.

\section*{Acknowledgements}
Computations used the PIK high-performance computing cluster. We thank Maik Billing for the global LPJmL-FIT
simulations used in Sect.~\ref{sec:global}, and the LPJmL-FIT developers for the model.

\bibliographystyle{plainnat}
\bibliography{refs}

\end{document}
